Die automatisierte, referenzfreie Bewertung von Sprachkomplexität bildet eine zentrale Grundlage dieser Arbeit. Um Texte in Leichter Sprache objektiv zu bewerten und ein stufenloses Belohnungssignal (\textit{Reward}) für das spätere Modelltraining bereitzustellen, wird in diesem Themenblock eine eigene Metrik-Pipeline entwickelt und systematisch evaluiert.

\subsection{Vorstufe: Binäre Klassifikation als Machbarkeitsnachweis \& Architektur-Baseline}
\label{sec:klassifikationsmodelle}

Bevor ein generatives Sprachmodell über überwachtes Feintuning (SFT) oder Präferenzoptimierung (DPO) zur Textvereinfachung trainiert werden kann, wird ein Bewertungsmodell benötigt, das den Grad sprachlicher Einfachheit objektiv und reproduzierbar misst.

\subsubsection{Motivation, Lesbarkeitsmetriken \& BiLSTM-Architektur}
\label{sec:klassifikation_motivation}

Klassische lesbarkeitsmetrische Ansätze wie der \textit{Flesch Reading Ease} \citep{amstad1978verstaendlich} oder die \textit{Wiener Sachtextformel} \citep{bamberger1984lesen} basieren primär auf statistischen Oberflächenmerkmalen wie der mittleren Silbenanzahl oder Satzlänge (vgl. Abschnitt~\ref{sec:linguistische_masse_lesbarkeit}). Zwar bieten diese Zählformeln eine recheneffiziente Basis zur groben Abschätzung allgemeiner Textschwierigkeit, sie stoßen bei der Bewertung von Leichter Sprache jedoch an Grenzen: Da sie für allgemeine Lesbarkeitsstufen im Bildungswesen entwickelt wurden, berücksichtigen sie die spezifischen Regeln der Leichten Sprache nicht. Weder syntaktische Vorgaben wie die Vermeidung von Passivkonstruktionen und Genitiven, die Auflösung von Nebensätzen in kurze Hauptsätze noch Besonderheiten wie Bindestrich-Kopplungen bei Komposita (z.\,B. \enquote{\textit{Werkstatt-Vertrag}}) werden von reinen Zählformeln erfasst. Zudem liefern feste Formeln keine stufenlosen Gradienten für das spätere Reward-Modeling im Präferenztraining.

Dieser Nachteil lässt sich durch einen datengetriebenen Ansatz beheben: Neuronale Modelle können die syntaktischen, morphologischen und lexikalischen Muster Leichter Sprache direkt aus den bereinigten Parallelkorpora lernen. Die binäre Klassifikation dient dabei als Machbarkeitsnachweis (\textit{Proof of Concept}): \textit{Ist ein rekurrentes Modell in der Lage, anhand der linguistischen Merkmale trennscharf zwischen alltagssprachlichen Ausgangstexten (AS, $y=0$) und Versionen in Leichter Sprache (LS, $y=1$) zu unterscheiden?}

Als neuronales Modell dient das in Abschnitt~\ref{sec:architekturen_sprachmodelle} theoretisch hergeleitete Bidirektionale LSTM (BiLSTM). Die Eingabetexte werden mittels \texttt{spaCy} tokenisiert und über eine Einbettungsschicht in zwei bidirektionale LSTM-Schichten eingespeist. Die aggregierte Sequenzrepräsentation wird nach Regularisierung über Dropout durch einen linearen Klassifikationskopf mit Sigmoid-Aktivierung in eine kontinuierliche Klassifikationswahrscheinlichkeit für Leichte Sprache überführt. Das Training erfolgt mittels binärer Kreuzentropie und dem AdamW-Optimierer \citep{loshchilov2019decoupled} mit Early Stopping. Die vollständigen Schichtdimensionen, Lernraten und Trainingsparameter aller Metrikmodelle sind in der Hyperparameter-Tabelle~\ref{tab:hyperparameter_metrik_modelle} im Anhang~\ref{sec:anhang_hyperparameter} dokumentiert.

Auf dem internen Test-Split des Master-Korpus bestätigt das Modell den ermittelten optimalen Ähnlichkeitsbereich ($0{,}80 \le \text{sim}_{\cos} \le 0{,}98$): Das BiLSTM erzielt eine Balanced Accuracy von $92{,}99\,\%$ auf Satzebene sowie $99{,}03\,\%$ auf Artikelebene (der vollständige In-Domain-Vergleich über Ähnlichkeitsbereiche ist in Tabelle~\ref{tab:sweet_spot_ablation_klassifikation} im Anhang~\ref{sec:anhang_klassifikatoren} dokumentiert). Dies belegt die hohe Qualität der Trainingsdaten und die Eignung des BiLSTM-Modells zur Erkennung einfacher Sprachstrukturen.

\subsubsection{Granularitäts-Vergleich: Satzebene vs. Artikelebene}
\label{sec:klassifikator_laengen_ablation}

Zur Bestimmung der optimalen Modellierungs- und Aggregationsebene sowie zur Überprüfung der Out-of-Domain-Generalisierung wurde eine systematische Vergleichsstudie auf dem unabhängigen, verifizierten \textit{Lebenshilfe}-Goldstandard ($N = 37$ Artikelpaare / $74$ Dokumente, $6.901$ Einzelsätze aus behördlich-juristischen Kontexten) durchgeführt. Untersucht wurden das Satzmodell (isoliert sowie aggregiert via Majority Voting) und direkte Artikel-Klassifikatoren über maximale Kontextlängen von $256$, $512$ und $1.024$ Tokens. Tabelle~\ref{tab:klassifikator_laengen_benchmark} stellt die Ergebnisse gegenüber klassischen Lesbarkeitsmetriken dar (ergänzende KDE-, ROC- und Konfusionsanalysen finden sich im Anhang~\ref{sec:anhang_klassifikatoren}).

\begin{table}[htbp]
\centering
\small
\caption[Benchmark der Klassifikationsmodelle auf dem Lebenshilfe-Goldstandard]{Benchmark der auf bidirektionalen Long-Short-Term-Memory-Netzwerken basierenden Klassifikationsmodelle (BiLSTM) über Granularitätsebenen und Sequenzlängen auf dem Lebenshilfe-Goldstandard ($N = 37$ Artikelpaare / $74$ Dokumente) im Vergleich zwischen Alltagssprache (AS) und Leichter Sprache (LS) inklusive Majority Voting über Einzelsätze (MajVote) und des relativen Paar-Rankings (Pair Match).}
\label{tab:klassifikator_laengen_benchmark}
\adjustbox{max width=\textwidth}{%
\begin{tabular}{@{}llccccc@{}}
\toprule
\textbf{Modell} & \textbf{Ebene / Input} & \textbf{\O{} Score (LS)} & \textbf{\O{} Score (AS)} & \textbf{Separation ($\Delta$)} & \textbf{Balanced Acc} & \textbf{Pair Match} \\
\midrule
BiLSTM Satz-Klassifikator          & Satzebene ($N=6.901$)   & $0{,}721 \pm 0{,}37$ & $0{,}160 \pm 0{,}29$ & $0{,}561$ & $79{,}94\,\%$ & -- \\
BiLSTM Satz-Klassifikator (MajVote)& Dokument ($r_{LS}$ Ratio)& $0{,}719 \pm 0{,}09$ & $0{,}116 \pm 0{,}13$ & $0{,}603$ & $\mathbf{98{,}65\,\%}$ & $\mathbf{100{,}0\,\%}\; (37/37)$ \\
\midrule
BiLSTM Artikel-Klassifikator ($256$)  & Dokument ($256$ Tok.)   & $0{,}756 \pm 0{,}36$ & $0{,}003 \pm 0{,}00$ & $0{,}754$ & $90{,}54\,\%$ & $97{,}3\,\%\; (36/37)$ \\
BiLSTM Artikel-Klassifikator ($512$)  & Dokument ($512$ Tok.)   & $0{,}781 \pm 0{,}36$ & $0{,}009 \pm 0{,}03$ & $0{,}772$ & $90{,}54\,\%$ & $94{,}6\,\%\; (35/37)$ \\
BiLSTM Artikel-Klassifikator ($1.024$) & Dokument ($1.024$ Tok.) & $0{,}413 \pm 0{,}40$ & $0{,}000 \pm 0{,}00$ & $0{,}413$ & $70{,}27\,\%$ & $\mathbf{100{,}0\,\%}\; (37/37)$ \\
\midrule
Flesch Reading Ease (Baseline)     & Dokument (Silben/Satz)  & $0{,}973 \pm 0{,}16$ & $0{,}189 \pm 0{,}39$ & $0{,}784$ & $89{,}19\,\%$ & $\mathbf{100{,}0\,\%}\; (37/37)$ \\
Wiener Sachtextformel (Baseline)   & Dokument (Wortlänge)    & $0{,}919 \pm 0{,}27$ & $0{,}108 \pm 0{,}31$ & $0{,}811$ & $90{,}54\,\%$ & $\mathbf{100{,}0\,\%}\; (37/37)$ \\
\bottomrule
\end{tabular}}
\end{table}

Die Ergebnisse liefern zwei fundamentale Erkenntnisse:
\begin{enumerate}
    \item \textbf{Satzmodell und Aggregation via Majority Voting:} Auf Einzelsatzebene fällt die Performanz mit $79{,}94\,\%$ Balanced Accuracy moderat aus, da isolierte Phrasen (z.\,B. kurze Überschriften oder Begrüßungen) oft keine eindeutigen Merkmale Leichter Sprache enthalten. Wird das Modell jedoch über die Simple-Sentence-Ratio
    \begin{equation}
        r_{LS} = \frac{1}{|S|} \sum_{s \in S} \mathbb{I}(\hat{y}_s = 1) \in [0{,}0, 1{,}0]
    \end{equation}
    mittels Majority Voting aggregiert, gleicht die Mittelung lokale Unsicherheiten weitgehend aus: Die Balanced Accuracy steigt auf $98{,}65\,\%$, der ROC-AUC auf $0{,}9949$ und das Paar-Matching erreicht fehlerfreie $100{,}0\,\%\; (37/37)$.
    
    \item \textbf{Signalverlust bei langen Artikeln:} Im direkten Kontrast dazu zeigen sich direkte Artikel-Klassifikatoren bei langen Sequenzen anfällig für Informationsverlust: Während bei $256$ und $512$ Tokens noch eine Balanced Accuracy von $90{,}54\,\%$ erzielt wird, bricht die Performanz bei $1.024$ Tokens auf $70{,}27\,\%$ ein. Die KDE-Analyse (Abbildung~\ref{fig:anhang_classifier_kde}) deckt eine deutliche Asymmetrie auf: Während Alltagssprache mit spitzen Peaks bei $0{,}0$ sehr sicher erkannt wird (\O{}-Score $\le 0{,}009$), überdecken neutrale Passagen und Eigennamen die Merkmale Leichter Sprache über lange Sequenzen hinweg. Der mittlere LS-Score sinkt auf $0{,}413$ ab und unterschreitet die Entscheidungsschwelle $\tau = 0{,}50$, wodurch der LS-Recall auf $40{,}54\,\%$ fällt (mehr falsch-negative Einstufungen).
\end{enumerate}

\paragraph{Einordnung der traditionellen Baselines gegenüber den neuronalen Modellen}
Die empirischen Ergebnisse in Tabelle~\ref{tab:klassifikator_laengen_benchmark} verdeutlichen die Stärken und Grenzen der traditionellen Formeln: Mit einer Balanced Accuracy von $89{,}19\,\%$ (Flesch) bzw. $90{,}54\,\%$ (Wiener Sachtextformel) und einem fehlerfreien Paar-Matching von $100{,}0\,\%$ ($37/37$) erweisen sich die klassischen Formeln als robuste Baselines auf Dokumentebene. Da Texte in Leichter Sprache deutlich kürzere Sätze und Wörter aufweisen, wirkt sich dieser Oberflächeneffekt auf Dokumentebene stark aus. Die Formeln übertreffen damit den direkten $1.024$-Token-Artikel-Klassifikator ($70{,}27\,\%$, $40{,}54\,\%$ LS-Recall), dessen Leistung durch die Signalüberlagerung über lange Sequenzen und die geringe Zahl eigenständiger Trainingsartikel abfällt.
Gegenüber dem hierarchisch aggregierten BiLSTM-Satzmodell ($r_{LS}$) mit $98{,}65\,\%$ Balanced Accuracy und einem ROC-AUC von $0{,}9949$ offenbaren die Formeln jedoch eine deutliche Schwäche: Ihre Dichteverteilungen (Abbildung~\ref{fig:anhang_classifier_kde}) weisen eine breite Überlappungszone im Übergangsbereich auf, da sie grammatikalische und stilistische Unterschiede nicht erfassen. Das neuronale Satzmodell schließt diese Lücke und erreicht eine trennscharfe Klassifikation ($\Delta = 0{,}603$).
\subsubsection{Dummy-Content-Test: Prüfung auf Längen- und Padding-Bias}
\label{sec:dummy_content_test}

Um auszuschließen, dass das BiLSTM lediglich oberflächliche Längenunterschiede oder das Zero-Padding (\textit{Shortcut Learning}) lernt, wurde ein gezielter Dummy-Content-Test auf dem Lebenshilfe-Goldstandard durchgeführt. Hierbei werden alle Wörter der Testdokumente durch ein einheitliches, neutrales Token (den Punkt \enquote{\texttt{.}} bzw. den Artikel \enquote{\textit{der}}) ersetzt, während die originären Sequenzlängen und das Zero-Padding exakt unverändert bleiben ($\mathbf{T}_{\text{dummy}} = (t_{\text{dummy}}, \dots, t_{\text{dummy}}, \text{PAD}, \dots)$). In diesem Szenario bricht die Balanced Accuracy von über $90\,\%$ auf exakt $50{,}00\,\%$ (reines Zufallsniveau) zusammen. Dies belegt, dass das Modell keine rein längen- oder paddingbasierten Heuristiken nutzt, sondern tatsächliche linguistische Merkmale verarbeitet.


\subsubsection{Synthese: Erlernbarkeit des Signals \& Überleitung zur stetigen Sprachstufenmodellierung}
\label{sec:grenzen_klassifikation_uebergang}

Zusammenfassend belegen die Ergebnisse, dass das BiLSTM linguistische Merkmale Leichter Sprache zuverlässig lernt und Dokumentenpaare verlässlich relativ rankt ($\text{ROC-AUC} > 0{,}94$, Pair-Match bis zu $100{,}0\,\%$). Für die spätere Steuerung generativer Modelle (Themenblock~3) erfordert das Fehlen kontinuierlicher Gradienten (Sigmoid-Sättigung des binären Klassifikationskopfes) jedoch ein stufenloses, differenzierbares Gütesignal $\lambda \in [0{,}0, 1{,}0]$.

Da natürliche Parallelkorpora der Leichten Sprache ausschließlich bimodal strukturiert sind (Ausgangstext in Alltagssprache vs. finale Übersetzung in Leichter Sprache) und in der Praxis keine authentischen Zwischenstufen vorliegen, muss ein kontinuierlicher Übergangsbereich modelliert werden. Im folgenden Abschnitt wird hierfür zunächst ein intuitiver Ansatz untersucht: die diskrete Generierung von Sprachkomplexitätsstufen mittels großer Sprachmodelle (LLMs).


\subsection{Erster Lösungsansatz: Synthetische Sprachstufengenerierung via Large Language Models}
\label{sec:synthetische_sprachstufen}

Als erster methodischer Ansatz zur Konstruktion eines stetigen Zwischenraums wurde die gezielte Generierung diskreter Sprachkomplexitätsstufen mittels großer Sprachmodelle (Large Language Models, LLMs) untersucht. Die Kernidee dieses Ansatzes bestand darin, ein Gütesignal auf Basis linguistisch modulierter Ganztexte zu erlernen, indem ein leistungsfähiges generatives Modell angewiesen wird, Zwischenversionen mit kontrolliertem Vereinfachungsgrad zu verfassen.

\subsubsection{Konzeption, Prompt-Design \& Generierung mit FlensGen-GPT}
\label{sec:synthetische_generierung}

Zur semantischen Generierung der Zwischenstufen wurde das an der Hochschule Flensburg betriebene Open-Source-Großsprachmodell \textbf{\texttt{FlensGen-GPT-OSS-120B}} über einen lokalen Inferenzendpunkt per Instruction Prompting angesteuert. Für jedes parallele Artikelpaar des bereinigten Master-Korpus sowie des Lebenshilfe-Datensatzes wurden unter Vorlage beider Originalfassungen (Alltagssprache als Stufe $0{,}00$ und Leichte Sprache als Stufe $1{,}00$) drei diskrete Zwischenstufen synthetisiert:
\begin{itemize}
    \item \textbf{Stufe $0{,}25$ (Nahe Alltagssprache):} Leicht vereinfachter Fließtext; Schachtelsätze werden aufgeteilt, schwierige Fachtermini kurz paraphrasiert.
    \item \textbf{Stufe $0{,}50$ (Einfache Sprache):} Typische \enquote{Einfache Sprache}; Verzicht auf Schachtelsätze bei Standard-Kombination aus Haupt- und Nebensatz, alltäglicher Wortschatz.
    \item \textbf{Stufe $0{,}75$ (Nahe Leichte Sprache):} Sehr leicht verständliche Sprache; überwiegend einfache Hauptsätze, jedoch ohne die extremen zeilenweisen Umbrüche Leichter Sprache.
\end{itemize}

Das vollständige System- und User-Prompt-Design inklusive aller Regelvorgaben ist in Anhang~\ref{sec:anhang_synthetic_regressor_prompt} dokumentiert. Auf dem resultierenden $5$-Stufen-Datensatz ($4.300$ Text-Samples aus $860$ Artikeln) wurde eine BiLSTM-Architektur identischer Kapazität (Hidden-Dim $128$, Embedding-Dim $128$, MSE-Loss) als \textit{Synthetischer Regressor} trainiert.

\subsubsection{Empirische Ergebnisse auf dem Lebenshilfe-Goldstandard}
\label{sec:synthetische_ergebnisse_vergleich}

Bei der empirischen Überprüfung auf dem unabhängigen \textit{Lebenshilfe}-Goldstandard zeigte der synthetische Regressor prinzipiell das erwünschte Verhalten: Das Modell erlernt die grundlegende Richtung der Sprachkomplexität und weist Texten in Leichter Sprache im Durchschnitt höhere Werte zu als den alltagssprachlichen Originalen ($\text{\O{}-Score}(\text{LS}) = 0{,}575 \pm 0{,}35$ vs. $\text{\O{}-Score}(\text{AS}) = 0{,}145 \pm 0{,}18$, $\Delta = +0{,}429$). Auch bei der relativen Gegenüberstellung gegebener Artikelpaare ordnet das Modell in $89{,}2\,\%$ der Fälle der vereinfachten Fassung einen höheren Score zu ($\text{ROC-AUC} = 0{,}8495$, Balanced Accuracy $78{,}38\,\%$, Pair Match $89{,}2\,\%\; [33/37]$). Entlang der fünf diskret generierten Stufen spiegelt sich zudem ein erkennbarer, stetiger Anstieg der vorhergesagten Werte wider.

Gleichzeitig offenbart die empirische Verteilungsanalyse jedoch erhebliche Schwächen (vgl. Abbildung~\ref{fig:anhang_synthetic_kde_roc} sowie Tabelle~\ref{tab:anhang_synthetic_vs_mixup_benchmark} im Anhang): Die Vorhersagen des synthetischen Regressors streuen spürbar breit, und insbesondere im mittleren Komplexitätsbereich fallen die Übergänge zwischen den Zwischenstufen diffus und überlappend aus.

\subsubsection{Methodische Grenzen \& Motivation für ein stochastisches Text-MixUp}
\label{sec:synthetische_grenzen_ausmusterung}

Obwohl der synthetische Ansatz als Machbarkeitsnachweis für LLM-basierte Komplexitätsmodulation funktionierte, erwies er sich für den praktischen Einsatz als Reward-Modell im generativen Alignment als unzureichend. Die Ursachen für diese Limitationen liegen in drei methodischen Kernproblemen begründet:
\begin{enumerate}
    \item \textbf{Granularitätsgrenze generativer Sprachmodelle:} Während das LLM grobe Niveaus trennen kann, stößt es bei feineren Zwischenstufen an eine Auflösungsgrenze. Eine Erhöhung auf beispielsweise $10$ oder $20$ Stufen erzeugt keinen Informationsgewinn, da aktuelle Modelle solch minimale Komplexitätsunterschiede nicht deterministisch steuern können und zu stilistischer Regression neigen.
    \item \textbf{Datenvolumen und statische Limitation:} Der synthetische Ansatz ist auf $5$ diskrete Texte pro Artikel limitiert ($4.300$ statische Samples). Bei Millionen Parametern neigt das rekurrente Netz bei wiederholter Epochen-Präsentation identischer Texte zu Overfitting.
    \item \textbf{Miterlernen künstlicher Modellierungsartefakte:} Der synthetische Ansatz trainiert auf maschinell generierten Zwischentexten. Dadurch lernt das rekurrente Netz unvermeidbare LLM-Artefakte, Halluzinationsmuster und unnatürliche Satzstrukturen als Scheinsignale mit, anstatt die authentischen sprachlichen Gesetzmäßigkeiten Leichter Sprache zu generalisieren.
\end{enumerate}

Aus diesen Gründen reicht die diskrete synthetische Generierung allein nicht aus, um ein stufenloses, robustes Belohnungssignal bereitzustellen. Im folgenden Abschnitt wird daher ein zweiter Lösungsansatz entwickelt: Ein stochastisches Mischverfahren, das ausschließlich auf authentischen, von Menschen erstellten Paralleltexten aufbaut und kontinuierliche Zielwerte erzeugt.


\subsection{Zweiter Lösungsansatz: Continuous MixUp Regression}
\label{sec:mixup_regression}

Um die Limitationen der synthetischen LLM-Generierung (Datenknappheit, diskrete Stufen und künstliche Textartefakte) zu überwinden, wird in diesem Abschnitt das Verfahren der \textit{Continuous MixUp Regression} entwickelt. Die Kernüberlegung besteht darin, den kontinuierlichen Übergangsraum direkt auf den authentischen, von menschlichen Fachexperten erstellten Parallelkorpora zu konstruieren.

\subsubsection{Konzeption des textuellen MixUp \& Dynamische Target-Berechnung}
\label{sec:mixup_konzeption}

Das klassische MixUp-Verfahren \citep{zhang2018mixup} erzeugt neue Trainingsdaten, indem es bestehende Eingaben und deren Zielwerte stufenlos miteinander vermischt. Da Wörter und Texte diskrete Strukturen sind und sich nicht direkt numerisch überblenden lassen, wird dieses Prinzip im Rahmen dieser Arbeit auf Satzebene übertragen: Kontinuierliche Zwischenstufen der Sprachkomplexität entstehen durch das kontrollierte, zufällige Mischen von Sätzen aus beiden Sprachfassungen:

Für jedes parallele Artikelpaar $(D_{\text{AS}}, D_{\text{LS}})$ des Trainingskorpus werden die Sätze segmentiert. Ein gemischtes Dokument $D_{\text{mix}}$ entsteht durch zufälliges Ziehen und Zusammenfügen von Sätzen aus beiden Sprachfassungen:
\begin{equation}
    D_{\text{mix}} = \text{Concat}\left( \text{sample}(S_{\text{LS}}, k_{\text{LS}}) \cup \text{sample}(S_{\text{AS}}, k_{\text{AS}}) \right)
\end{equation}
wobei die Sätze vor dem Zusammenfügen zufällig durchmischt werden, um zu verhindern, dass das rekurrente Netz bestimmte Sprachmuster rein anhand ihrer Position im Text lernt.

Da Sätze in Alltagssprache und Leichter Sprache stark divergierende Wort- und Satzlängen aufweisen, würde ein reiner Satzanzahl-Anteil den tatsächlichen sprachlichen Anteil systematisch verzerren. Der kontinuierliche Zielwert $\lambda \in [0{,}0, 1{,}0]$ wird daher dynamisch über den relativen Zeichenlängenanteil der Leichten Sprache im gemischten Text berechnet:
\begin{equation}
    \lambda = \frac{\text{CharLen}(S_{\text{LS, gezogen}})}{\text{CharLen}(S_{\text{LS, gezogen}}) + \text{CharLen}(S_{\text{AS, gezogen}})}
\end{equation}
Hierbei repräsentiert $\lambda = 0{,}0$ reine Alltagssprache und $\lambda = 1{,}0$ reine Leichte Sprache, während alle Zwischenwerte stufenlose Simplizitätsgrade quantifizieren.

\subsubsection{Vergleich der Trainings- und Sampling-Strategien}
Bei der methodischen Entwicklung des textuellen MixUp-Trainings wurden vier Dataloader- und Optimierungsstrategien systematisch verglichen:
\begin{itemize}
    \item \textbf{Variante A (Statisch):} Feste Vorabmischung des Trainingskorpus bei der Dataset-Initialisierung ($p_{\text{dynamic}} = 0{,}0$).
    \item \textbf{Variante B (Dynamisch):} Rein dynamisches On-the-Fly-Shuffling in jedem Batch ($p_{\text{dynamic}} = 1{,}0$).
    \item \textbf{Variante C (Hybrid):} Epochenabhängiges Curriculum-Mischen mit kontinuierlichem Übergang von statischen zu dynamischen Mischungen ($p_{\text{dynamic}} = \frac{\text{Epoche}}{\text{Gesamtepochen}}$).
    \item \textbf{Variante D (Hybrid + Cyclic LR):} Hybride Mischung kombiniert mit zyklischem Lernratenscheduling mittels \textit{CosineAnnealingWarmRestarts} nach \citet{loshchilov2016sgdr} zur Vermeidung lokaler Minima.
\end{itemize}

\paragraph{Einfluss der Trainings- und Sampling-Strategien auf die Regressionsgüte}
In frühen Experimenten mit stark begrenzter Datenmenge ($M = 10$ Mischungen pro Artikelpaar, entsprechend ca.~$5.130$ Samples) zeigten sich deutliche Unterschiede zwischen den Ansätzen: Die statische Vorabmischung neigte durch die wiederholte Präsentation identischer Satzkombinationen zu Overfitting. Die rein dynamische Variante konvergierte bei wenigen Iterationen ($N=513$) aufgrund des hohen Gradientenrauschens nur langsam. Die hybride Variante mit zyklischem Scheduler erwies sich in dieser Phase als am stärksten, da sie eine stabile anfängliche Repräsentationsbildung ermöglichte und durch periodische Lernratenanstiege lokalen Minima entkommen konnte.

Wird der Trainingskorpus jedoch voll skaliert ($M = 160$, über $82.000$ Mischungen pro Epoche über $100$ Epochen), gleichen sich die Ergebnisse der vier Ansätze weitgehend an: Durch die große Zahl unterschiedlicher Satzkombinationen konvergieren alle vier Varianten zu einem nahezu identischen Leistungsniveau ($\text{Val-MSE} \approx 0{,}027$, Test-Balanced-Accuracy $\approx 96\,\%$, vgl. Tabelle~\ref{tab:anhang_mixup_varianten_benchmark} im Anhang~\ref{sec:anhang_mixup_varianten}). Die Wahl der Dataloader-Strategie ist somit vor allem in datenlimitierten Szenarien ausschlaggebend, während bei umfassender Datenaugmentation die Modellkapazität des BiLSTMs die Obergrenze vorgibt. Die detaillierten Dichteverteilungen und Regressionsanalysen aller vier Varianten sind in Abbildung~\ref{fig:anhang_mixup_variants_kde} und Abbildung~\ref{fig:anhang_mixup_variants_scatter} im Anhang dokumentiert.

\subsubsection{Einfluss der Kontextlänge auf die MixUp-Regression}
\label{sec:regressor_laengen_ablation}

Analog zur Klassifikation wurde der Einfluss der Eingabesequenzlänge auf das kontinuierliche MixUp-Regressionsmodell über $256$, $512$ und $1.024$ Tokens systematisch evaluiert. Das Training erfolgte auf über $80.000$ generierten dynamischen MixUp-Samples mit kontinuierlichen Zielwerten $\lambda \in [0{,}0, 1{,}0]$.

Bereits während der Validierung auf dem Test-Split zeigte sich ein kontinuierlicher Vorteil längerer Kontexte: Der minimale Mean Squared Error (MSE) sank von $0{,}0260$ ($256$ Tokens) über $0{,}0246$ ($512$ Tokens) auf $0{,}0234$ bei $1.024$ Tokens (minimaler MAE $0{,}0976$). Tabelle~\ref{tab:regressor_laengen_benchmark} stellt die Evaluierung der Regressormodelle auf dem externen Lebenshilfe-Goldstandard dar (ergänzende grafische Analysen der Dichteverteilungen finden sich in Abbildung~\ref{fig:anhang_regressor_kde} im Anhang).

\begin{table}[htbp]
\centering
\small
\caption[Längen-Benchmark der BiLSTM MixUp Regressoren auf dem Lebenshilfe-Goldstandard]{Längen-Benchmark der auf bidirektionalen Long-Short-Term-Memory-Netzwerken basierenden MixUp-Regressoren (BiLSTM) auf dem Lebenshilfe-Goldstandard ($N = 37$ Artikelpaare / $74$ Dokumente) zur kontinuierlichen Einstufung der Textmischung $\lambda \in [0{,}0, 1{,}0]$ zwischen Alltagssprache (AS, $\lambda=0{,}0$) und Leichter Sprache (LS, $\lambda=1{,}0$) inklusive ausgewogener Genauigkeit (Balanced Accuracy bei $\tau=0{,}50$), ROC-AUC und relativem Paar-Ranking (Pair Match).}
\label{tab:regressor_laengen_benchmark}
\adjustbox{max width=\textwidth}{%
\begin{tabular}{@{}llcccccc@{}}
\toprule
\textbf{Modell} & \textbf{Ebene / Input} & \textbf{\O{} Score (LS)} & \textbf{\O{} Score (AS)} & \textbf{Separation ($\Delta$)} & \textbf{Balanced Acc} & \textbf{ROC-AUC} & \textbf{Pair Match} \\
\midrule
BiLSTM MixUp Regressor ($256$)  & Dokument ($256$ Tok.)   & $0{,}698 \pm 0{,}11$ & $0{,}105 \pm 0{,}09$ & $0{,}593$ & $95{,}95\,\%$ & $0{,}9982$ & $97{,}3\,\%\; (36/37)$ \\
BiLSTM MixUp Regressor ($512$)  & Dokument ($512$ Tok.)   & $0{,}716 \pm 0{,}12$ & $\mathbf{0{,}069 \pm 0{,}09}$ & $\mathbf{0{,}646}$ & $\mathbf{98{,}65\,\%}$ & $0{,}9993$ & $97{,}3\,\%\; (36/37)$ \\
BiLSTM MixUp Regressor ($1.024$) & Dokument ($1.024$ Tok.) & $\mathbf{0{,}727 \pm 0{,}14}$ & $0{,}100 \pm 0{,}09$ & $0{,}627$ & $97{,}30\,\%$ & $\mathbf{1{,}0000}$ & $\mathbf{100{,}0\,\%}\; (37/37)$ \\
\midrule
Flesch Reading Ease (Baseline)     & Dokument (Silben/Satz)  & $0{,}973 \pm 0{,}16$ & $0{,}189 \pm 0{,}39$ & $0{,}784$ & $89{,}19\,\%$ & $0{,}8919$ & $100{,}0\,\%\; (37/37)$ \\
Wiener Sachtextformel (Baseline)   & Dokument (Wortlänge)    & $0{,}919 \pm 0{,}27$ & $0{,}108 \pm 0{,}31$ & $0{,}811$ & $90{,}54\,\%$ & $0{,}9054$ & $100{,}0\,\%\; (37/37)$ \\
\bottomrule
\end{tabular}}
\end{table}

\paragraph{Inverses Skalierungsverhalten: Warum dominiert beim Regressor das $1.024$-Token-Modell?}
Im direkten Kontrast zum Artikel-Klassifikator erweist sich bei der MixUp-Regression das $1.024$-Token-Modell als das stärkste und stabilste Modell:
\begin{enumerate}
    \item \textbf{Perfektes Paar-Ranking ($100{,}0\,\%$):} Bei ausnahmslos allen $37$ von $37$ Artikelpaaren des Goldstandards ordnet das $1.024$-Modell der Leichten Sprache einen strikt höheren Score zu als der zugehörigen Alltagssprache ($\text{Score}(\text{LS}) > \text{Score}(\text{AS})$, vgl. Tabelle~\ref{tab:regressor_laengen_benchmark}).
    \item \textbf{Höchste Trennschärfe \& perfekter ROC-AUC:} Das Modell erreicht den höchsten mittleren LS-Score ($0{,}727$) und erzielt ein fehlerfreies Paar-Ranking mit einem perfekten ROC-AUC von $\mathbf{1{,}0000}$.
\end{enumerate}

\paragraph{Systematische Überlegenheit des MixUp-Regressors gegenüber Klassifikator und synthetischem Regressor}
Die direkte Gegenüberstellung der Benchmark-Ergebnisse (Tabelle~\ref{tab:klassifikator_laengen_benchmark} und Tabelle~\ref{tab:regressor_laengen_benchmark}) liefert eine zentrale methodische Erkenntnis dieser Arbeit: Während das direkte Artikel-Training beim binären Klassifikator bei langen Sequenzen instabil wird und der synthetische Regressor bei $\text{ROC-AUC} = 0{,}8495$ stagniert, bleibt das Training des BiLSTM MixUp-Regressors hochgradig robust, skaliert positiv mit längeren Kontexten und erzielt auf dem Goldstandard einen fehlerfreien ROC-AUC von $\mathbf{1{,}0000}$.

Die Ursachen für diese Überlegenheit begründen sich in vier Mechanismen:
\begin{itemize}
    \item \textbf{Gleichmäßige Merkmalsverteilung durch Satz-Mischung:} Beim Artikel-Klassifikator stehen die Sätze in ihrer festen Originalreihenfolge. Treten Merkmale Leichter Sprache nur am Textanfang auf und folgt danach langer Fließtext, verliert das LSTM über viele Rekursionsschritte das Signal. Beim MixUp-Regressor werden die gezogenen Sätze vor dem Zusammenfügen zufällig durchmischt. Dadurch verteilen sich Merkmale beider Sprachstufen gleichmäßig über das gesamte $1.024$-Token-Fenster (Anfang, Mitte und Ende), was einem positionsabhängigen Signalverlust vorbeugt.
    \item \textbf{Kontinuierliche Datenvielfalt statt Datenmangel:} Auf Dokumentebene stehen dem Klassifikator und dem synthetischen Ansatz nur wenige feste Trainingsartikel zur Verfügung ($750$ bzw. $860$ Artikel), was das Risiko von Overfitting erhöht. Der MixUp-Regressor generiert hingegen in jeder Epoche zehntausende neu gemischte Satzkombinationen und lernt somit auf einer kontinuierlich variierenden Datenbasis.
    \item \textbf{MSE-Verlustfunktion statt binärer Entscheidung:} Die binäre Klassifikation (BCE-Loss) neigt dazu, Ausgaben schnell auf $0$ oder $1$ festzulegen. Schon einzelne schwere Wörter können zu Fehlentscheidungen führen. Die MSE-Verlustfunktion bestraft Abweichungen dagegen gleichmäßig über das gesamte Intervall $[0{,}0, 1{,}0]$ und sorgt für stabile, stufenlose Vorhersagen.
    \item \textbf{Erfassung des relativen Textanteils:} Während der Klassifikator eine harte Entscheidung zwischen $0$ und $1$ treffen muss, misst der Regressor den relativen Anteil einfacher Sätze im Text. Ein größeres Kontextfenster ($1.024$ Tokens) bietet dabei eine breitere Stichprobe und senkt die Vorhersageschwankungen auf $\pm 0{,}10$.
\end{itemize}

\paragraph{MixUp-Regression vs. Traditionelle Formel-Baselines}
Gegenüber den klassischen Formeln erzielt der $1.024$-Token MixUp-Regressor eine deutlich höhere Trennschärfe ($97{,}30\,\%$ Balanced Accuracy, ROC-AUC $1{,}0000$). Traditionelle Formeln sind an feste Schwellenwerte gebunden, streuen im mittleren Bereich stark und geben keine stufenlosen Werte aus. Das MixUp-Modell liefert dagegen kontinuierliche und differenzierbare Scores. Dadurch kann im späteren Modelltraining auch ein kleiner Fortschritt bei der Textvereinfachung präzise belohnt werden.


\subsection{Validierung der kontinuierlichen Abstufungsfähigkeit auf externen Benchmarks (TextComplexityDE \& MERLIN)}
\label{sec:validierung_abstufung_externe_benchmarks}

Die vorangegangene Evaluierung auf dem \textit{Lebenshilfe}-Goldstandard (Abschnitt~\ref{sec:regressor_laengen_ablation}) belegt, dass der BiLSTM MixUp-Regressor ein trennscharfes Modell zur Unterscheidung auf Dokumentebene ist. Trotz des kontinuierlichen MixUp-Trainingsverfahrens unterscheidet das Modell mit einer Balanced Accuracy von $97{,}30\,\%$, einem ROC-AUC von $1{,}0000$ und einem fehlerfreien Perfect-Pair-Matching von $100{,}0\,\%$ ($37/37$) zwischen Alltagssprache und Leichter Sprache und übertrifft sowohl den originär binär trainierten Klassifikator als auch den synthetischen Regressor systematisch.

\subsubsection{Das Problem fehlender Zwischenstufenkorpora \& Motivation der externen Benchmarks}
\label{sec:dilemma_zwischenstufen}

Die vorangegangene Evaluierung auf dem \textit{Lebenshilfe}-Goldstandard belegt die hervorragende Diskriminationskraft des MixUp-Regressors auf den beiden Extrempolen (Alltagssprache vs. Leichte Sprache). Das eigentliche Ziel der Regressionsmodellierung besteht jedoch darin, ein stufenloses Gütesignal bereitzustellen, das den kontinuierlichen Grad sprachlicher Simplizität quantifiziert.

Die direkte empirische Überprüfung dieser kontinuierlichen Regressionsfähigkeit auf realen Texten stößt auf eine methodische Hürde: In der Praxis existieren keine natürlichen Parallelkorpora, die Zwischenstufen zwischen Alltagssprache und Leichter Sprache enthalten. Die professionelle Übersetzung in Leichte Sprache erfolgt prinzipbedingt zweistufig (komplexer Ausgangstext vs. maximal vereinfachte Zielübersetzung). Reale Übergangszustände (beispielsweise ein Text mit einfachen Sätzen, aber noch langen Komposita) sind in keinem Korpus systematisch erfasst.

Um die Regressionsgüte des Modells dennoch unabhängig von den künstlichen Trainingsmischungen zu evaluieren, stellen externe linguistische Datensätze einen methodischen Versuch dar, die kontinuierliche Einstufungsfähigkeit im Zero-Shot-Transfer zu überprüfen. Zwar bilden auch diese Benchmarks keine direkte Interpolation zwischen Alltagssprache und den normierten Regeln Leichter Sprache ab, sie bieten jedoch standardisierte, unabhängige Komplexitätsskalen:
\begin{enumerate}
    \item \textbf{TextComplexityDE \citep{naderi2019textcomplexityde}:} Ein standardisierter Benchmark aus $1.000$ isolierten deutschen Wikipedia-Sätzen mit aggregierten menschlichen Urteilen (Likert-Skala) zur wahrgenommenen Satzschwierigkeit und Verständlichkeit.
    \item \textbf{MERLIN-Korpus \citep{boyd2014merlin}:} Ein Korpus aus $1.033$ authentischen deutschen Volltexten auf Dokumentebene, die von Sprachexperten offiziell nach den sechs Niveaustufen des Gemeinsamen Europäischen Referenzrahmens für Sprachen (CEFR: A1, A2, B1, B2, C1, C2) eingestuft wurden.
\end{enumerate}

\subsubsection{Satzebene: Evaluation auf TextComplexityDE \& Analyse der Längendynamik}
\label{sec:eval_textcomplexityde}

Der \textit{TextComplexityDE}-Benchmark \citep{naderi2019textcomplexityde} umfasst $1.000$ deutsche Wikipedia-Sätze, für die aggregierte menschliche Urteile (\textit{Mean Opinion Score}, MOS) hinsichtlich Gesamtschwierigkeit ($\text{MOS}_{\text{Complexity}} \in [1{,}0, 7{,}0]$), Verständlichkeit und lexikalischer Hürden vorliegen. Während ein in früheren Auswerteskripten genutztes statisches Zero-Padding die Hidden-States bei kurzen Einzelsätzen durch hunderte Padding-Schritte stark verwässerte und die Ergebnisse massiv verschlechterte, ermöglichte die Umstellung auf dynamisches Sequenz-Padding eine unverfälschte Inferenz mit deutlich präziseren Vorhersagen. Die Untersuchung in Abbildung~\ref{fig:textcomplexityde_boxplots} liefert wesentliche methodische Erkenntnisse über das Modellverhalten auf Satzebene:
\begin{itemize}
    \item \textbf{Überlegenheit des $256$-Token-Modells auf Satzebene:} Die visuelle und statistische Analyse der Monotonie-Boxplots (Abbildung~\ref{fig:textcomplexityde_boxplots}) offenbart einen zentralen methodischen Befund: Der $256$-Token-Regressor ist auf Satzebene das am saubersten abgestufte Modell. Mit dem geringsten normierten Fehler ($\text{MAE} = 0{,}198$) und der höchsten linearen Korrelation ($r = 0{,}620$) nutzt das Modell die Skala mit gleichmäßig fallenden Medianen optimal aus.
    \item \textbf{Stauchung der Werte bei langen Texten:} Während das $256$-Modell einen Einzelsatz als gewichtigen Anteil seines Trainingsfensters interpretiert, entspricht ein isolierter Satz für das $1.024$-Modell nur rund $1{,}5\,\%$ des Kontextes. Enthält ein Wikipedia-Satz auch nur ein einzelnes Fach- oder Standardwort (z.\,B. Eigennamen oder Institutionen), ordnet das $1.024$-Token-Modell den kurzen Textausschnitt primär der Alltagssprache zu. Dadurch werden die Vorhersagen des $1.024$-Token-Modells bei isolierten Einzelsätzen ab Stufe 2 in einen schmalen unteren Wertebereich verschoben.
    \item \textbf{Kontrast Satz- vs. Dokumentebene:} Dieser Befund begründet eine klare methodische Trennung: Auf isolierten Einzelsätzen dominiert das $256$-Token-Modell in Skalennutzung und Dynamik, während auf vollständigen Artikeln (Lebenshilfe-Goldstandard, $100\,\%$ Perfect Pair Match) das $1.024$-Token-Modell überlegen ist, da es lange Dokumente ohne Informationsverlust durch Truncation ganzheitlich abbildet.
\end{itemize}

Aufgrund dieser methodischen Besonderheiten bei isolierten Einzelsätzen ist die ergänzende Evaluation auf zusammenhängenden Fließtexten (MERLIN) für eine ganzheitliche Beurteilung unerlässlich.

\begin{figure}[H]
\centering
\includegraphics[width=\textwidth]{images/textcomplexityde_boxplots_all_models_grid.png}
\caption[Verteilungs-Boxplots über 5 Likert-Komplexitätsstufen auf TextComplexityDE]{Verteilungs-Boxplots der Vorhersagewerte neuronaler BiLSTM-Modelle über fünf diskrete menschliche Likert-Komplexitätsstufen auf dem TextComplexityDE-Benchmark ($N = 1.000$ Sätze, Stufe 1 bis 5) mit Darstellung des Interquartilsabstands (IQR).}
\label{fig:textcomplexityde_boxplots}
\end{figure}

\subsubsection{Dokumentebene: Evaluation auf dem MERLIN CEFR-Benchmark}
\label{sec:eval_merlin}

Um die kontinuierliche Sprachstufung auf echten Fließtexten ohne Längen- und Padding-Verzerrungen zu evaluieren, wird das MERLIN-Korpus \citep{boyd2014merlin} herangezogen. Mit $1.033$ vollständigen Texten (Aufsätze und Briefe aus standardisierten Sprachprüfungen) und einer durchschnittlichen Textlänge von ca. $150$ bis $250$ Tokens entspricht dieses Korpus exakt der typischen Dokumentstruktur von Leichte-Sprache-Absätzen.

Tabelle~\ref{tab:anhang_merlin_master} und Tabelle~\ref{tab:anhang_merlin_monotonie} im Anhang~\ref{sec:anhang_merlin} dokumentieren die detaillierten Kennwerte aller Modelle über die CEFR-Stufen. Die Auswertung in Abbildung~\ref{fig:merlin_boxplots} liefert wesentliche Erkenntnisse über die graduelle Einstufungsfähigkeit:
\begin{itemize}
    \item \textbf{Monotoner Verlauf des $256$-Token MixUp-Regressors:} Wie die Verteilungen in Abbildung~\ref{fig:merlin_boxplots} verdeutlichen, sinken die vorhergesagten Simplizitätsscores des $256$-Token-Regressors über die sechs CEFR-Niveaustufen von der leichtesten Stufe A1 bis zur Muttersprachenstufe C2 kontinuierlich und streng monoton ab. Dies zeigt, dass das Modell auf vollständigen Texten eine konsistente Sprachstufung vornimmt ($\rho = 0{,}485$, Pearson $r = 0{,}459$, $\text{MAE} = 0{,}310$).
    \item \textbf{Ausgeprägte Niveautrennung der $512$er- und $1.024$er-MixUp-Regressoren:} Auch die Modelle mit größerem Kontextfenster bilden die CEFR-Schwierigkeit ab: Während elementare Sprachstufen (A1, A2, B1) im Bereich von $0{,}24$ bis $0{,}34$ liegen, vollzieht sich ab dem Übergang zur fortgeschrittenen Sprachverwendung (B2: $\sim 0{,}16$--$0{,}19$) sowie zu Fach- und Muttersprachenniveau (C1/C2: $\sim 0{,}08$--$0{,}12$) ein stetiger, hochgradig signifikanter Score-Abfall (Tabelle~\ref{tab:anhang_merlin_monotonie}).
    \item \textbf{Korrelation der Baselines mit CEFR-Prüfungsstufen:} Auch auf MERLIN korrelieren traditionelle Formeln stark mit den Niveaus (LIX: $\rho = 0{,}656$, WSTF: $\rho = 0{,}602$, Flesch: $\rho = 0{,}543$), da höhere Prüfungsstufen (B2--C2) formal anspruchsvollere Wort- und Satzstrukturen fordern. Wie zuvor diskutiert, bildet dies eine nützliche Basis, vermag jedoch die spezifischen, normierten Vereinfachungsregeln Leichter Sprache nicht feinkörnig zu erfassen oder als Gradient für generative Modelle bereitzustellen.
\end{itemize}

\begin{figure}[H]
\centering
\includegraphics[width=\textwidth]{images/merlin_boxplots_all_models_grid.png}
\caption[Monotonie-Prüfung über CEFR-Sprachniveaustufen auf MERLIN]{Monotonie-Prüfung der Regressoren und Baselines über die Sprachniveaustufen des Gemeinsamen Europäischen Referenzrahmens (CEFR, A1 bis C2) auf dem MERLIN-Benchmark ($N = 1.033$ Texte) anhand von Boxplots mit Kennzeichnung von Interquartilsabstand, Median und Mittelwert.}
\label{fig:merlin_boxplots}
\end{figure}

\subsubsection{Kritische Synthese: Validität des Zero-Shot-Transfers \& Fazit}
\label{sec:synthese_abstufung_fazit}

Bei der Interpretation dieses Benchmark-Transfers ist eine methodische Differenzierung geboten: Weder TextComplexityDE noch das MERLIN-Korpus wurden für Leichte Sprache konzipiert. Das CEFR-Raster klassifiziert die Fremdsprachenkompetenz von Deutschlernenden (L2-Spracherwerb) und sanktioniert orthografische sowie grammatikalische Fehler, während Leichte Sprache festen Regeln folgt (u.\,a. nach BITV 2.0 \citep{bitv2019} und den Richtlinien des BMAS \citep{bmas2025ratgeber}). Diese fordern vor allem kurze Hauptsätze, einfache Wörter und getrennte Wortzusammensetzungen.

Gerade vor diesem Hintergrund ist das empirische Abschneiden der BiLSTM-MixUp-Regressoren als bemerkenswert zu bewerten: Obwohl die Modelle originär ausschließlich darauf trainiert wurden, die Mischungsanteile zwischen Alltagssprache und Leichter Sprache stufenlos einzuschätzen, behaupten sie sich auf diesen gänzlich domänenfremden (Out-of-Domain) und auf völlig andere Sprachdimensionen abzielenden Korpora hervorragend. Auf isolierten Einzelsätzen (TextComplexityDE) übertrifft der $256$-Token-Regressor traditionelle Formeln in der linearen Skalennutzung ($r = 0{,}620$, $\text{MAE} = 0{,}198$), während auf Dokumentebene (MERLIN) ein streng monotoner, stabiler Score-Abfall über alle sechs CEFR-Stufen ($0{,}371 \to 0{,}078$) nachgewiesen werden konnte.

Dies liefert den entscheidenden empirischen Nachweis, dass das MixUp-Verfahren keine isolierten Datensatz-Artefakte erzeugt hat. Stattdessen lernt das neuronale Netz ein verlässliches, übertragbares Verständnis sprachlicher Einfachheit, das sich im weiteren Verlauf dieser Arbeit verlässlich als stetiges Belohnungssignal ($R_{\text{style}}$) für das anschließende Modell-Alignment (Themenblock~3) eignet.

\subsection{Synthese \& Modellauswahl für die generative Übersetzung (Themenblock 3)}
\label{sec:synthese_modellauswahl_themenblock3}

Um die empirische Güte und Trennschärfe aller in diesem Themenblock entwickelten Ansätze vergleichend zusammenzufassen, stellt Tabelle~\ref{tab:master_benchmark_themenblock2} die Ergebnisse sämtlicher Modellfamilien auf dem unabhängigen \textit{Lebenshilfe}-Goldstandard ($N = 37$ Artikelpaare / $74$ Dokumente) gegenüber.

\paragraph{Vergleich mit einem einfachen RNN (Vanilla RNN)}
Um zu überprüfen, ob die bidirektionale LSTM-Zellstruktur mit ihren Gating-Mechanismen (Input-, Forget- und Output-Gates, vgl. Abschnitt~\ref{sec:architekturen_sprachmodelle}) für die Komplexitätsmodellierung notwendig ist, wurde ein einfaches unidirektionales Netz (\textit{Vanilla RNN} nach Elman) als Vergleichsmodell unter identischen Bedingungen trainiert. Ohne Gating-Mechanismen leidet das einfache RNN bei längeren Sequenzen unter verschwindenden Gradienten (\textit{Vanishing Gradient Problem}). Zwar erreicht die Baseline ein relatives Ranking ($\text{Balanced Accuracy} = 67{,}57\,\%$, $\text{ROC-AUC} = 0{,}7761$, Pair Match $70{,}3\,\%$, $\text{Score}(\text{LS}) = 0{,}443$, $\text{Score}(\text{AS}) = 0{,}345$), ihre Trennschärfe fällt mit $\Delta = +0{,}098$ jedoch deutlich geringer aus als beim BiLSTM-Modell ($\Delta = +0{,}627$). Dies belegt, dass die Gates des BiLSTM notwendig sind, um längere Texte zuverlässig zu verarbeiten.

\begin{table}[htbp]
\centering
\small
\caption[Gesamtübersicht aller in Themenblock 2 evaluierten Modellansätze]{Gesamtübersicht aller in Themenblock 2 evaluierten Modellansätze auf dem Lebenshilfe-Goldstandard ($N = 37$ Artikelpaare / $74$ Dokumente) zur kontinuierlichen Einstufung oder Klassifikation zwischen Alltagssprache (AS) und Leichter Sprache (LS) unter Einbeziehung des Text-MixUp-Verfahrens, großer Sprachmodelle (Large Language Models, LLM) und rekurrenten Basisarchitekturen (Recurrent Neural Networks, RNN).}
\label{tab:master_benchmark_themenblock2}
\adjustbox{max width=\textwidth}{%
\begin{tabular}{@{}llcccccc@{}}
\toprule
\textbf{Modell} & \textbf{Paradigma / Input} & \textbf{\O{} Score (LS)} & \textbf{\O{} Score (AS)} & \textbf{Separation ($\Delta$)} & \textbf{Balanced Acc} & \textbf{ROC-AUC} & \textbf{Pair Match} \\
\midrule
\multicolumn{8}{@{}l}{\textit{Kontinuierliche Regressionsmodelle (MixUp)}} \\
BiLSTM MixUp Regressor ($1.024$) & Dokument ($1.024$ Tok.) & $\mathbf{0{,}727 \pm 0{,}14}$ & $0{,}100 \pm 0{,}09$ & $0{,}627$ & $97{,}30\,\%$ & $\mathbf{1{,}0000}$ & $\mathbf{100{,}0\,\%}\; (37/37)$ \\
BiLSTM MixUp Regressor ($512$)  & Dokument ($512$ Tok.)   & $0{,}716 \pm 0{,}12$ & $\mathbf{0{,}069 \pm 0{,}09}$ & $\mathbf{0{,}646}$ & $\mathbf{98{,}65\,\%}$ & $0{,}9993$ & $97{,}3\,\%\; (36/37)$ \\
BiLSTM MixUp Regressor ($256$)  & Dokument ($256$ Tok.)   & $0{,}698 \pm 0{,}11$ & $0{,}105 \pm 0{,}09$ & $0{,}593$ & $95{,}95\,\%$ & $0{,}9982$ & $97{,}3\,\%\; (36/37)$ \\
\midrule
\multicolumn{8}{@{}l}{\textit{Diskrete \& Synthetische Regressionsmodelle}} \\
BiLSTM Synthetischer Regressor  & Dokument (LLM-Stufen)   & $0{,}575 \pm 0{,}35$ & $0{,}145 \pm 0{,}18$ & $0{,}429$ & $78{,}38\,\%$ & $0{,}8495$ & $89{,}2\,\%\; (33/37)$ \\
Vanilla RNN Baseline (Vergleich) & Dokument ($256$ Tok.)   & $0{,}443 \pm 0{,}20$ & $0{,}345 \pm 0{,}18$ & $0{,}098$ & $67{,}57\,\%$ & $0{,}7761$ & $70{,}3\,\%\; (26/37)$ \\
\midrule
\multicolumn{8}{@{}l}{\textit{Binäre Klassifikationsmodelle}} \\
BiLSTM Satz-Klassifikator (MajVote)& Dokument ($r_{LS}$ Ratio)& $0{,}769 \pm 0{,}08$ & $0{,}127 \pm 0{,}14$ & $0{,}642$ & $\mathbf{98{,}65\,\%}$ & $\mathbf{0{,}9985}$ & $\mathbf{100{,}0\,\%}\; (37/37)$ \\
BiLSTM Artikel-Klassifikator ($256$)& Dokument ($256$ Tok.)   & $0{,}829 \pm 0{,}31$ & $\mathbf{0{,}002 \pm 0{,}01}$ & $\mathbf{0{,}827}$ & $91{,}89\,\%$ & $0{,}9974$ & $97{,}3\,\%\; (36/37)$ \\
BiLSTM Artikel-Klassifikator ($512$)& Dokument ($512$ Tok.)   & $0{,}478 \pm 0{,}41$ & $0{,}031 \pm 0{,}13$ & $0{,}446$ & $71{,}62\,\%$ & $0{,}8809$ & $91{,}9\,\%\; (34/37)$ \\
BiLSTM Artikel-Klassifikator ($1.024$)& Dokument ($1.024$ Tok.)& $0{,}340 \pm 0{,}40$ & $0{,}007 \pm 0{,}01$ & $0{,}333$ & $67{,}57\,\%$ & $0{,}8795$ & $97{,}3\,\%\; (36/37)$ \\
\midrule
\multicolumn{8}{@{}l}{\textit{Traditionelle Formeln (Baselines)}} \\
Flesch Reading Ease (Baseline)  & Dokument (Silben/Satz)  & $0{,}973 \pm 0{,}16$ & $0{,}189 \pm 0{,}39$ & $0{,}784$ & $89{,}19\,\%$ & $0{,}8919$ & $100{,}0\,\%\; (37/37)$ \\
Wiener Sachtextformel (Baseline)& Dokument (Wortlänge)    & $0{,}919 \pm 0{,}27$ & $0{,}108 \pm 0{,}31$ & $0{,}811$ & $90{,}54\,\%$ & $0{,}9054$ & $100{,}0\,\%\; (37/37)$ \\
\bottomrule
\end{tabular}}
\end{table}

\paragraph{Synthese: Scheinbare Überlegenheit der Baselines vs. Die Teilmengen-Hypothese}
Auf den ersten Blick vermitteln die Ergebnisse in Tabelle~\ref{tab:master_benchmark_themenblock2} den Eindruck, als seien traditionelle Zählformeln wie der Flesch Reading Ease ($89{,}19\,\%$ BAcc, $100{,}0\,\%$ Pair Match) oder die Wiener Sachtextformel ($90{,}54\,\%$ BAcc, $100{,}0\,\%$ Pair Match) die einfachere, robustere und damit vorzuziehende Wahl zur Simplizitätsbewertung.

Dieser Eindruck täuscht jedoch: Die makroskopische Treffsicherheit traditioneller Formeln auf dem \textit{Lebenshilfe}-Goldstandard beruht ausschließlich darauf, dass menschliche Fachübersetzer in Leichter Sprache alle Regeln gleichzeitig umsetzen. Klassische Formeln messen davon jedoch lediglich eine oberflächliche Teilmenge ($S_{\text{metric}} \subset S_{\text{LS}}$), namentlich die mittlere Satzlänge ($ASL$) und die Silbendichte ($ASW$). Da authentische Leichte Sprache extrem kurze Sätze und elementare Wörter nutzt, deckt sie diese Teilmenge zwangsläufig mit ab.

Sobald Texte jedoch nicht mehr perfekt vereinfacht sind (wie es bei fehlerhaften Modellgenerierungen im Übersetzungsprozess typischerweise vorkommt), versagen traditionelle Formeln. Dies belegt ein kontrollierter synthetischer Benchmark ($N = 30$ Textfassungen aus fünf Themenbereichen mit jeweils $100$ bis $256$ Tokens), bei dem systematische grammatikalische Regelverletzungen isoliert auf Dokumentebene injiziert wurden (die vollständige Ergebnistabelle und exemplarische Textfassungen sind in Anhang~\ref{sec:anhang_synthetic_rule_benchmark} sowie Tabelle~\ref{tab:anhang_synthetic_rule_benchmark_table} dokumentiert):
\begin{itemize}
    \item \textbf{Blindheit gegenüber dem Konjunktiv ($\Delta = 0{,}000$):} Werden alle Verben eines Leichte-Sprache-Textes konsequent vom Indikativ in den unzulässigen Konjunktiv II überführt (\enquote{\textit{haben}} $\rightarrow$ \enquote{\textit{hätten}}, \enquote{\textit{kann}} $\rightarrow$ \enquote{\textit{könnte}}, \enquote{\textit{ist}} $\rightarrow$ \enquote{\textit{wäre}}), bleibt die Silben- und Tokenanzahl exakt gleich ($\Delta \text{Silben} = 0$). Für den Flesch Reading Ease und die WSTF-4 ist ein unzulässiger Konjunktiv-Text von korrekter Leichter Sprache nicht zu unterscheiden ($\Delta = 0{,}000$).
    \item \textbf{Fehlerhafte Belohnung von Genitivkonstruktionen ($+6{,}55$ Flesch-Punkte):} Werden regelkonforme Umschreibungen durch Genitive ersetzt (\enquote{\textit{das Büro von dem Arzt}} $\rightarrow$ \enquote{\textit{das Büro des Arztes}} oder \enquote{\textit{der Schlüssel für das Haus}} $\rightarrow$ \enquote{\textit{der Schlüssel des Hauses}}), entfallen Wörter (\enquote{\textit{von dem}} $\rightarrow$ \enquote{\textit{des}}). Flesch registriert die kürzere Satzlänge ($ASL$) und stuft den fehlerhaften Text fälschlicherweise als einfacher ein ($+6{,}55$ Punkte Anstieg von $54{,}25$ auf $60{,}81$).
    \item \textbf{Fehlerhafte Belohnung von Passivstrukturen ($+6{,}04$ Flesch-Punkte):} Vorgangspassive fügen Hilfsverben ein (\enquote{\textit{wird von}}), welche die durchschnittliche Silbenanzahl pro Wort rechnerisch senken. Flesch belohnt diesen Regelverstoß fälschlicherweise mit $+6{,}04$ Punkten Verbesserung ($54{,}25 \rightarrow 60{,}29$).
\end{itemize}

Im direkten Kontrast dazu reagiert der trainierte $256$-Token BiLSTM-Regressor empfindlich auf diese Regelbrüche: Bei Konjunktiv-Injektion fällt der Modell-Score um $-62{,}4\,\%$ und bei Genitiv-Injektion um $-54{,}1\,\%$ ab. Dies belegt empirisch, dass das neuronale Metrikmodell grammatikalische Regeln erfasst, die über reine Zählstatistiken hinausgehen. Zusammen mit der Bereitstellung stufenloser Werte qualifiziert dieses Regelverständnis das neuronale Modell als valides Belohnungssignal ($R_{\text{style}}$) für die nachfolgende Präferenzoptimierung.

\paragraph{Finale Modellauswahl für Themenblock 3}
Auf Basis dieses empirischen Gesamtergebnisses wird für die nachfolgenden Übersetzungs- und Alignment-Stufen (Supervised Fine-Tuning [SFT], Abschnitt~\ref{sec:sft_training}, sowie Direct Preference Optimization [DPO], Abschnitt~\ref{sec:dpo_optimization}) verbindlich das $1.024$-Token BiLSTM MixUp-Regressionsmodell als primäre Einfachheitsmetrik ($R_{\text{style}}$) ausgewählt.

Diese Entscheidung stützt sich auf vier zentrale Faktoren:
\begin{enumerate}
    \item \textbf{Empirische Überlegenheit in den Metrik-Tests:} Im Längenvergleich auf dem unabhängigen \textit{Lebenshilfe}-Goldstandard (Abschnitt~\ref{sec:regressor_laengen_ablation}, Tabelle~\ref{tab:regressor_laengen_benchmark}) erzielte die $1.024$-Token-Variante das beste Gesamtergebnis mit einem fehlerfreien Paar-Matching von $100{,}0\,\%$ ($37/37$ Paare), maximaler Separation, fehlerfreiem ROC-AUC ($1{,}0000$) sowie dem niedrigsten Validierungs-MSE ($0{,}0234$).
    \item \textbf{Systematische Überlegenheit gegenüber dem synthetischen Regressor:} Gegenüber dem LLM-basierten Synthetischen Regressor ($\text{BAcc} = 78{,}38\,\%$, $\text{ROC-AUC} = 0{,}8495$, Pair Match $89{,}2\,\%$, Abschnitt~\ref{sec:synthetische_sprachstufen}) bietet das MixUp-Modell eine deutlich höhere Trennschärfe und Robustheit auf realen Texten.
    \item \textbf{Maximale Artikelabdeckung:} Im bereinigten Korpus besitzen Ausgangstexte im Mittel $622{,}2$ Wörter ($\approx 1.350$ BPE-Tokens) und vereinfachte Texte $457{,}9$ Wörter ($\approx 857$ BPE-Tokens). Da rund $45\,\%$ der Artikel $512$ Tokens und über $90\,\%$ das $128$-Token-Limit überschreiten, gewährleistet erst das $1.024$-Token-Fenster eine ganzheitliche Dokumentenbewertung ohne vorzeitiges Abschneiden (\textit{Truncation}).
    \item \textbf{Gleiche Kontextlänge wie beim Übersetzungsmodell (mBART-50):} Das als Übersetzungsmodell eingesetzte Sequenz-zu-Sequenz-Modell \texttt{facebook/mbart-large-50} \citep{tang2020multilingual, liu2020multilingual} besitzt eine native Positionskodierungsgrenze von maximal $1.024$ Tokens (\texttt{max\_position\_embeddings = 1024}). Die Angleichung der Metrik auf $1.024$ Tokens stellt sicher, dass das Belohnungssignal im DPO-Training exakt das vollständige Generierungsfenster von mBART-50 abdeckt.
\end{enumerate}
